% Standalone problem document, generated from the adjacent README.md.
% Compile with XeLaTeX. Mathematical content is maintained in Markdown.
\documentclass[11pt,a4paper]{article}
\usepackage[fontsize=10.5pt]{scrextend}
\usepackage[margin=22mm,headheight=15pt,headsep=9mm,footskip=11mm]{geometry}
\usepackage{fontspec}
\setmainfont{texgyrepagella}[Extension=.otf,UprightFont=*-regular,BoldFont=*-bold,ItalicFont=*-italic,BoldItalicFont=*-bolditalic]
\setsansfont{texgyreheros}[Extension=.otf,UprightFont=*-regular,BoldFont=*-bold,ItalicFont=*-italic,BoldItalicFont=*-bolditalic]
\setmonofont{texgyrecursor}[Extension=.otf,UprightFont=*-regular,BoldFont=*-bold,ItalicFont=*-italic,BoldItalicFont=*-bolditalic,Scale=MatchLowercase]
\usepackage{amsmath,amssymb}
\usepackage{unicode-math}
\setmathfont{texgyrepagella-math.otf}
\usepackage{xcolor,microtype,enumitem,needspace,fancyhdr}
\usepackage{bookmark}
\definecolor{ink}{HTML}{17344B}
\definecolor{accent}{HTML}{197D80}
\definecolor{muted}{HTML}{596773}
\hypersetup{colorlinks=true,linkcolor=accent,urlcolor=accent,pdftitle={SP-12 - A
chromatic lower bound for positive-semidefinite nullity with
SAP},pdfauthor={Open Problems in Numerical Linear Algebra}}
\urlstyle{same}
\setlength{\parindent}{0pt}
\setlength{\parskip}{5pt plus 1pt minus 1pt}
\linespread{1.04}
\setlength{\emergencystretch}{3em}
\widowpenalty=10000
\clubpenalty=10000
\displaywidowpenalty=10000
\allowdisplaybreaks[1]
\setlist{nosep,leftmargin=1.5em}
\providecommand{\tightlist}{\setlength{\itemsep}{0pt}\setlength{\parskip}{0pt}}
\providecommand{\pandocbounded}[1]{#1}
\pagestyle{fancy}
\fancyhf{}
\fancyhead[L]{\sffamily\footnotesize\color{muted}OPEN PROBLEMS IN NUMERICAL LINEAR ALGEBRA}
\fancyhead[R]{\sffamily\bfseries\footnotesize\color{accent}SP-12}
\fancyfoot[L]{\parbox[b]{0.9\textwidth}{\sffamily\fontsize{8}{10}\selectfont\color{muted}Editorial ratings; a bounded search cannot certify that a problem remains unsolved.\\Literature check: 2026-09-11}}
\fancyfoot[R]{\sffamily\footnotesize\color{muted}\thepage}
\renewcommand{\headrulewidth}{0.3pt}
\renewcommand{\headrule}{\hbox to\headwidth{\color{accent}\leaders\hrule height \headrulewidth\hfill}}
\makeatletter
\renewcommand{\section}{\Needspace{5\baselineskip}\@startsection{section}{1}{0pt}{12pt plus 2pt minus 2pt}{4pt}{\sffamily\bfseries\large\color{ink}}}
\renewcommand{\subsection}{\Needspace{4\baselineskip}\@startsection{subsection}{2}{0pt}{9pt plus 1pt minus 1pt}{3pt}{\sffamily\bfseries\normalsize\color{ink}}}
\makeatother
\setcounter{secnumdepth}{0}
\begin{document}
{\sffamily\small\color{muted}Eigenvalues and inverse problems\par}
\vspace{3pt}
{\raggedright\hyphenpenalty=10000\exhyphenpenalty=10000\sffamily\bfseries\fontsize{21}{24}\selectfont\color{ink}A
chromatic lower bound for positive-semidefinite nullity with SAP\par}
\vspace{7pt}
{\sffamily\small\textbf{Difficulty:} challenging\quad\textbf{Importance:} interesting
to specialist\par}
{\sffamily\small\bfseries\color{accent}Status: Solved\par}
\vspace{4pt}
{\color{accent}\hrule height 0.8pt}
\vspace{6pt}
\textbf{Rating rationale (historical):} Challenging reflects the need
for a stronger general PSD/SAP realization bound than current
connectivity estimates supply; implication from Hadwiger does not
establish equivalent difficulty. Specialist impact reflects the focus on
this particular graph nullity invariant.

\subsection{Literature-dependent resolution -
2026-09-11}\label{literature-dependent-resolution---2026-09-11}

\textbf{Solved.} Hall's \textbf{Corollary 3.22}, based on
\textbf{Theorem 3.20}, gives \(\nu(H)\ge\delta(H)\) for every graph
\(H\). Induced-subgraph monotonicity and a vertex-critical subgraph of
chromatic number \(\chi(G)\) then prove the exact target
\(\nu(G)\ge\chi(G)-1\). The application note supplies the complete
monotonicity and coloring argument, including rank-zero and disconnected
cases.

The all-graph theorem is due to \textbf{H. Tracy Hall},
\href{https://arxiv.org/html/2601.01211v1}{\emph{The Delta Theorem},
arXiv:2601.01211v1}, submitted 3 January 2026. \textbf{George
Stepaniants}, Department of Computing and Mathematical Sciences,
California Institute of Technology, is the author of the
\href{https://github.com/ajt60gaibb/OpenProblemsInNLA/blob/main/eigenvalues-and-inverse-problems/SP-12/solution.md}{explanatory
application note}; no new theorem discovery or priority is claimed.
\href{https://github.com/ajt60gaibb/OpenProblemsInNLA/blob/main/eigenvalues-and-inverse-problems/SP-12/solution.pdf}{Application
PDF} ·
\href{https://github.com/ajt60gaibb/OpenProblemsInNLA/blob/main/eigenvalues-and-inverse-problems/SP-12/solution.tex}{Standalone
TeX}.

A separate
\href{https://github.com/ajt60gaibb/OpenProblemsInNLA/blob/main/references/stepaniants-sp11-sp12-2026-09-11/verification/SP-11-SP-12-independent-review.md}{Codex-agent
review} checked the full essential proof in Hall's preprint and the
exact deduction and returned \textbf{PASS}. The source remains a
preprint; this is independent automated-agent review, not external human
peer review or formal verification. Substantial ChatGPT/Codex assistance
and the review's precise limits are disclosed in the
\href{https://github.com/ajt60gaibb/OpenProblemsInNLA/blob/main/references/stepaniants-sp11-sp12-2026-09-11/README.md}{submission
record}.

The earlier status checks below are retained as historical records. The
original ID, target, path, historical ratings and attributed partial
results remain unchanged. Unsupported prior artifact and graph-atlas
claims remain withdrawn.

\newpage
\subsection{Problem statement}\label{problem-statement}

Let \(G\) be a finite simple undirected graph on \(n\ge1\) vertices. Let
\(\mathcal S(G)\) be the real symmetric matrices whose off-diagonal
nonzero entries occur exactly at edges of \(G\), with unrestricted
diagonal.

A matrix \(A\in\mathcal S(G)\) has the \textbf{strong Arnold property
(SAP)} if the only real symmetric matrix \(X\) satisfying \[
AX=0,\qquad A\circ X=0,\qquad I_n\circ X=0
\] is \(X=0\), where \(\circ\) is the entrywise product. Define \[
\nu(G)=\max\{\dim\ker A:A\in\mathcal S(G),\ A\succeq0,\ A\text{ has SAP}\}.
\] Writing \(\chi(G)\) for the minimum number of colors in a proper
vertex coloring, does every \(G\) satisfy \[
\nu(G)\ge\chi(G)-1?
\]

\subsection{Relevance and ratings}\label{relevance-and-ratings}

The conclusion asserts the existence of a positive-semidefinite matrix,
with an exact prescribed sparsity pattern, of rank at most
\(n-\chi(G)+1\) and with an additional nondegeneracy property.
Equivalently it constrains structured Gram matrix construction. This is
adjacent to core NLA through matrix realization and inverse eigenvalue
theory.

\newpage
\subsection{References}\label{references}

\begin{itemize}
\item
  F. Barioli, S. M. Fallat, H. Gupta, and Z. Li, \emph{The weak version
  of the graph complement conjecture and partial results for the delta
  conjecture}, Discrete Mathematics 349 (2026), 114861, Conjecture 4.1
  and §1 definitions (\href{https://arxiv.org/html/2505.24577v1}{primary
  manuscript};
  \href{https://doi.org/10.1016/j.disc.2025.114861}{journal}).
\item
  F. Barioli, S. M. Fallat, and L. Hogben, \emph{A variant on the graph
  parameters of Colin de Verdière: Implications to the minimum rank of
  graphs}, Electronic Journal of Linear Algebra 13 (2005), 387--404,
  pp.388--390 on SAP and nullity parameters
  (\href{https://journals.uwyo.edu/index.php/ela/article/download/341/341}{primary
  paper}).
\item
  M. Chudnovsky and A. Ovetsky Fradkin, \emph{Hadwiger's conjecture for
  quasi-line graphs}, Journal of Graph Theory 59 (2008), 17--33, Theorem
  1.1
  (\href{https://web.math.princeton.edu/~mchudnov/Hadwiger_quasiline.pdf}{primary
  manuscript}; \href{https://doi.org/10.1002/jgt.20321}{journal}).
\end{itemize}

\subsection{Status check}\label{status-check}

This exact inequality is expressly proposed in the 2026 paper. Its
weaker independence-number consequence, Conjecture 4.4, is not
separately counted. The source explains that the strong delta conjecture
and Hadwiger's conjecture would each imply the assertion; the original
unrestricted symmetric delta conjecture in the companion entry is a
different statement. Searches for this paper, ``nu'', ``chromatic'',
``positive semidefinite'', ``conjecture'', and 2026 found no resolution.
No claim is made that these conjectures are logically independent.

The 2026 counterexamples to a Colin de Verdière conjecture about Laplace
eigenvalue multiplicities on surfaces concern a different invariant;
they do not refute this finite PSD/SAP inequality.

\subsection{Audit update --- 2026-09-10}\label{audit-update-2026-09-10}

\textbf{Partially resolved:} the inequality holds for quasi-line graphs,
meaning each vertex neighborhood is the union of two cliques. This
follows by combining Chudnovsky--Fradkin Theorem 1.1 with the proved
inequality \(\nu(G)\ge\eta(G)-1\) recalled immediately before Conjecture
4.1, where \(\eta\) is the largest clique-minor order. This is an
implication of two established results, not a claim that the general
conjecture is proved. Both full primary texts, the 2026 publication
record, and targeted later searches were checked; no general resolution
was located.
\end{document}
